% =====================================================================
% protocolo.tex
% Responsabilidad única: dibujar el protocolo temporal del proyecto
% (entrenamiento, validación y prueba por fechas; qué se hace con cada bloque;
% validación cruzada de ventana creciente y punto de congelamiento). Fechas y
% tamaños son macros generadas por los notebooks.
% =====================================================================
\begin{figure}[!htbp]
\centering
\begin{tikzpicture}[
    every node/.style={font=\small},
    caja/.style={draw, rounded corners=2pt, align=center, minimum height=1.05cm, inner sep=4pt},
    paso/.style={draw, rounded corners=2pt, align=center, fill=gray!6, text width=4.2cm,
                 minimum height=1.9cm, inner sep=4pt},
    flecha/.style={-{Stealth[length=2.2mm]}, thick},
]
% --- Datos y partición por fechas ------------------------------------------
\node[caja, fill=gray!10, text width=15.2cm] (datos) at (0,0)
    {\textbf{Datos:}\enspace \ExpMuestrasLimpias{} muestras limpias recibidas entre \DatDesde{} y \DatHasta{}\\[1pt]
     {\scriptsize\color{gray!70!black}partición por fecha de recepción, sin mezclar: Cortes en \DatCorteValidacion{} y \DatCortePrueba{}}};

\node[caja, fill=classblue!14, draw=classblue, text width=4.2cm] (tr) at (-5.3,-2.2)
    {\textbf{Entrenamiento}\\ antes de \DatCorteValidacion{}\\ \ExpFilasEntrenamiento{} muestras, \RegDiasEntrenamiento{} días};
\node[caja, fill=classgreen!16, draw=classgreen, text width=4.2cm] (val) at (0,-2.2)
    {\textbf{Validación}\\ \DatCorteValidacion{} a \DatCortePrueba{}\\ \ExpFilasValidacion{} muestras, \RegDiasValidacion{} días};
\node[caja, fill=classred!12, draw=classred, text width=4.2cm] (te) at (5.3,-2.2)
    {\textbf{Prueba}\\ desde \DatCortePrueba{}\\ \ExpFilasPrueba{} muestras, \FinDiasPrueba{} días};

\foreach \n in {tr, val, te} \draw[flecha, gray] (datos.south -| \n) -- (\n.north);

% --- Qué se hace con cada bloque ------------------------------------------
\node[paso, below=0.75cm of tr] (ajuste)
    {\textbf{Ajustar y explorar}\\ cribado, L16 de Taguchi y rejilla, con validación cruzada temporal};
\node[paso, below=0.75cm of val] (decidir)
    {\textbf{Decidir y congelar}\\ modelo, calibración y umbral por costo; Reajuste con entrenamiento + validación};
\node[paso, below=0.75cm of te] (evaluar)
    {\textbf{Evaluar una sola vez}\\ Etapa 07: Métricas con intervalos bootstrap, sin cambiar nada};

\draw[flecha, classblue] (tr) -- (ajuste);
\draw[flecha, classgreen] (val) -- (decidir);
\draw[flecha, classred] (te) -- (evaluar);
\draw[flecha] (ajuste) -- (decidir);
\draw[flecha] (decidir) -- (evaluar);

% --- Barrera de congelamiento ----------------------------------------------
\coordinate (barrera) at ($(te.west)+(-0.28,0)$);
\draw[classred, very thick, dashed] ($(barrera)+(0,0.9)$) -- ($(barrera |- evaluar.south)+(0,-0.1)$);
\node[classred, font=\scriptsize\bfseries, align=center, fill=white, inner sep=1pt]
    at ($(barrera)+(0,1.2)$) {lectura única y registrada};

% --- Validación cruzada de ventana creciente (dentro de entrenamiento) ------
\begin{scope}[shift={($(ajuste.south)+(-2.1,-0.55)$)}]
  \node[font=\scriptsize, anchor=west, gray!70!black] at (0,0)
      {\DatBloquesCV{} bloques de ventana creciente (solo entrenamiento):};
  \foreach \i/\a in {1/1.6, 2/2.3, 3/3.0, 4/3.7} {
    \fill[classblue!45] (0,-0.28*\i-0.12) rectangle (\a,-0.28*\i+0.04);
    \fill[classgreen!70] (\a,-0.28*\i-0.12) rectangle (\a+0.7,-0.28*\i+0.04);
  }
  \node[font=\scriptsize, anchor=west] at (4.6,-0.55) {\textcolor{classblue}{ajusta}};
  \node[font=\scriptsize, anchor=west] at (4.6,-0.9) {\textcolor{classgreen!70!black}{valida}};
\end{scope}
\end{tikzpicture}
\caption{Protocolo temporal del proyecto, donde cada bloque de fechas cumple
un solo papel: Con entrenamiento ajustamos y exploramos configuraciones con
validación cruzada de ventana creciente (abajo a la izquierda, donde cada
pliegue entrena con el pasado y valida con el tramo siguiente). Con
validación decidimos y congelamos. La prueba se abre una sola vez en la
etapa~07, que deja registrada cada lectura.}
\label{fig:protocolo}
\end{figure}
